\documentclass[twoside]{article}

\usepackage[preprint]{aistats2027}
\usepackage[round,authoryear]{natbib}
\usepackage{amsmath,amssymb,amsthm}
\usepackage{mathtools}
\usepackage{microtype}
\usepackage{booktabs}
\usepackage{url}

\renewcommand{\bibname}{References}
\renewcommand{\bibsection}{\subsubsection*{\bibname}}
\bibliographystyle{apalike}

\newtheorem{theorem}{Theorem}
\newtheorem{lemma}[theorem]{Lemma}
\newtheorem{proposition}[theorem]{Proposition}
\newtheorem{definition}[theorem]{Definition}
\newcommand{\E}{\mathbb E}
\newcommand{\R}{\mathbb R}
\newcommand{\cF}{\mathcal F}
\newcommand{\cO}{\mathcal O}
\newcommand{\ve}{\varepsilon}
\newcommand{\prox}{\operatorname{prox}}
\newcommand{\supp}{\operatorname{supp}}
\newcommand{\Span}{\operatorname{span}}

\begin{document}
\runningauthor{}

\twocolumn[
\aistatstitle{No Randomized Speedup for\\ Weakly Convex Optimization}

\aistatsauthor{Pahan Dewasurendra}

\aistatsaddress{Johns Hopkins University}
]

\begin{abstract}
For globally Lipschitz weakly convex optimization, the best known methods in
the local value-and-subdifferential model need order $\ve^{-4}$ queries to make
the Moreau-envelope gradient at most $\ve$. A recent lower bound proves that
this rate is unavoidable for
deterministic algorithms, even when every query returns the full
subdifferential set. Whether randomization can improve the rate was left open.

We prove that it cannot. Every randomized algorithm requires
\[
 \Omega\!\left(
 \frac{G^2\min\{\rho\Delta,G^2\}}{\ve^4}
 \right)
\]
queries on the class of $G$-Lipschitz, $\rho$-weakly convex functions with
initial gap at most $\Delta$. The oracle returns both the value and the entire
subdifferential. The hard instances live in a polynomial dimension.
When $\Delta\leq G^2/\rho$, the lower bound matches the known randomized
upper bound in all four problem parameters.

Two obstacles distinguish the result from a standard random rotation. A full
subdifferential exposes any exact dependence on an unrevealed direction, and
unbounded queries defeat concentration. We give the diagonal chain a smooth
dead zone that yields exact local independence. We then combine a smooth soft
projection with a linear-growth radial penalty. This wrapper is globally
Lipschitz, preserves weak convexity, hides unbounded queries, and introduces no
spurious Moreau-stationary point. A conditional-Haar argument completes the
randomized lower bound.
\end{abstract}

\section{Introduction}

A function $f:\R^d\to\R$ is $\rho$-weakly convex if
$f+\frac{\rho}{2}\|\cdot\|^2$ is convex. This class contains convex-composite,
robust-estimation, and phase-retrieval objectives while allowing nonsmooth and
nonconvex geometry. Its standard stationarity measure is the gradient of the
Moreau envelope
\[
 f_{1/(2\rho)}(x)=\min_z\{f(z)+\rho\|x-z\|^2\}.
\]
Subgradient and model-based methods find a point with
$\E\|\nabla f_{1/(2\rho)}(x)\|\leq\ve$ using
$O(\rho G^2\Delta/\ve^4)$ queries in the small-gap regime
\citep{davis2018stochastic,davis2019model}. The fourth power contrasts with
the second-power rate of smooth nonconvex optimization.
Earlier model-based work established asymptotic convergence without a rate
\citep{duchi2018stochastic}. More recent deterministic work gives
fourth-power guarantees for related Goldstein targets and for Moreau
stationarity through bundle subproblems
\citep{kong2024cost,liao2025proximal}.

\citet{li2026deterministic} recently showed that this contrast is intrinsic
for deterministic algorithms. Their oracle is unusually strong. At $x$, it
returns $(f(x),\partial f(x))$, including the entire subdifferential set rather
than a selected subgradient. Their construction couples tilted max-affine
blocks into a constant-width diagonal zero-chain. A finite-horizon resisting
rotation then handles arbitrary deterministic queries.

That rotation can depend on a deterministic trajectory. It does not apply to
a randomized method. Drawing a rotation in advance creates two leaks. First,
a small positive projection onto an unrevealed direction changes the smooth
relay and appears in the full subdifferential. Second, a randomized algorithm
can query at arbitrarily large norm, amplifying a random projection until the
hidden direction becomes visible. The latter issue also appears in smooth
nonconvex lower bounds, where it is handled by a soft projection and a
quadratic penalty \citep{carmon2020stationary}. A quadratic is not available
here because the objective must be globally Lipschitz.

We resolve both leaks. A shifted smoothstep makes the diagonal chain locally
constant on a neighborhood of every unrevealed relay coordinate. The tilted
facets already have strict intercept gaps. At radius
$\tau=1/(64N^2)$, these two facts give exact local independence from every
direction more than one filtration layer ahead. This exact statement matters:
an approximation in value or gradient would not control an oracle that returns
the full subdifferential set.

For unbounded queries, we introduce the wrapper
\begin{equation}
 H_U(x)=\frac14\left[
 \bar f\bigl(U^\top s_S(x)\bigr)
 +3\bigl(\sqrt{B^2+\|x\|^2}-B\bigr)
 \right],
 \label{eq:wrapper}
\end{equation}
where $s_S(x)=x/\sqrt{1+\|x\|^2/S^2}$ and $U$ has orthonormal columns.
The first term only sees a point of norm at most $S$. The second has bounded
slope and prevents stationarity at infinity. Their interaction is the main
technical issue. We prove that the radial term cannot cancel the chain
subgradient at a proximal point. Every positive coordinate on the support of
that subgradient is $O(1/N)$, while the chain norm is
$\Omega(N^{-1/2})$.

\paragraph{Contributions.}
Our main theorem closes the randomized local-oracle complexity in the small-gap
regime $\Delta\leq G^2/\rho$. It matches both the deterministic lower bound
and the known randomized upper bound, including the dependence on
$G,\rho,\Delta$, and $\ve$. The proof also gives a reusable globally Lipschitz
replacement for the quadratic soft-projection wrapper in randomized
stationarity lower bounds. Finally, we formulate a robust full-set chain
condition and show that a width-six diagonal filtration is compatible with the
conditional-Haar argument of \citet{carmon2020stationary}.

\paragraph{Related lower bounds.}
Random rotations are classical in randomized convex lower bounds
\citep{woodworth2017randomized,diakonikolas2020parallel}. The robust zero-chain
and soft-projection framework of \citet{carmon2020stationary} treats smooth
functions and local derivative oracles. \citet{kornowski2022oracle} prove a
different impossibility result for finding a point near one with a small
Clarke subgradient. Their class is unrestricted Lipschitz nonconvex functions,
and their target is not Moreau stationarity. The closest result is
\citet{li2026deterministic}, which proves the deterministic version of our rate
and explicitly leaves randomized full-set algorithms open. We retain their
tilted-block architecture, but both robustness and the Lipschitz wrapper are
new.

The upper-bound literature uses several oracle models. Stochastic
subgradient and model-based schemes control the Moreau-envelope gradient
\citep{davis2018stochastic,davis2019model}, while proximally guided methods
control a nearby stationarity certificate \citep{davis2019guided}. The
deterministic method of \citet{kong2024cost} uses a directional subgradient
oracle. \citet{liao2025proximal} obtain the same fourth-power Moreau rate with
function values and subgradients through convex bundle subproblems. Our lower
bound grants the algorithm the full set $\partial f(x)$, so it covers these
first-order information models when specialized to the same class and target.

Structured smoothing can be faster under stronger access. The framework of
\citet{deng2026smooth} assumes that a prescribed smooth approximation is
computationally tractable and that its gradients can be queried. Moreau
smoothing then requires solving a proximal optimization problem, while
Nesterov smoothing uses a known composite representation. The method of
\citet{najar2026variable} similarly uses explicit composite structure and a
linear minimization oracle. Such surrogate gradients are not determined by
the local pair $(f(x),\partial f(x))$. Our theorem therefore separates these
structured models from black-box local first-order access.

\section{Model and Main Result}
\label{sec:model}

For $G,\rho,\Delta>0$, let $\cF(G,\rho,\Delta)$ be the union over dimensions of finite-valued,
bounded-below functions $f:\R^d\to\R$ such that $f$ is $G$-Lipschitz,
$\rho$-weakly convex, and
\[
 f(0)-\inf_x f(x)\leq\Delta.
\]
The oracle is $\cO_f(x)=(f(x),\partial f(x))$. Weakly convex functions are
subdifferentially regular, so the usual limiting, Clarke, and regular
subdifferentials agree \citep{rockafellar2009variational}. Algorithms are
measurable, may use internal randomness, start at zero, and may choose each
query from the full previous transcript. An arbitrary transcript-measurable
output is allowed.

\begin{definition}[Randomized query complexity]
For failure probability $\delta\in(0,1)$, let
$T_{\ve,\delta}^{\rm rand}(\cF)$ be the least $T$ for which a randomized
algorithm can, for every $f\in\cF$, make at most $T$ oracle calls and output
$\widehat x$ satisfying
$\Pr(\|\nabla f_{1/(2\rho)}(\widehat x)\|\leq\ve)\geq1-\delta$.
\end{definition}

\begin{theorem}[Randomization does not help]
\label{thm:main}
There are numerical constants $c_0,c_1>0$ such that, whenever
\[
 \ve^2\leq c_0\min\{G^2,\rho\Delta\},
\]
we have
\[
 T_{\ve,1/2}^{\rm rand}\bigl(\cF(G,\rho,\Delta)\bigr)
 \geq c_1
 \frac{G^2\min\{\rho\Delta,G^2\}}{\ve^4}.
\]
The hard function can be taken in a dimension polynomial in $G/\ve$.
\end{theorem}

When $\Delta\leq G^2/\rho$, projected subgradient descent
\citep{davis2018stochastic} and Theorem~\ref{thm:main} give the minimax rate
\begin{equation}
 \Theta\!\left(
 \frac{\rho G^2\Delta}{\ve^4}
 \right).
 \label{eq:minimax-rate}
\end{equation}
Indeed, the optimized bound of \citet[Corollary~2.2]{davis2018stochastic},
followed by Markov's inequality, gives success probability at least $1/2$ at
this query order.
For larger gaps, Theorem~\ref{thm:main} still gives the saturated lower bound
$\Omega(G^4/\ve^4)$, but we do not claim a matching gap-independent upper
bound.

\section{An Exactly Robust Diagonal Chain}
\label{sec:chain}

We first construct an unscaled instance on $M$ blocks of $N$ coordinates.
Assume $N\geq60$ and $2\leq M\leq N/12$. Write
$z=(z^1,\ldots,z^M)\in\R^{MN}$. Define
\begin{align*}
 c_j&=1+\frac{N-j}{16N},\\
 \phi(u)&=\max\left\{0,\max_{j\in[N]}
 \left(c_j-\frac N8u_j\right)\right\}.
\end{align*}
Let $q(t)$ equal zero for $t\leq0$, $3t^2-2t^3$ for $0<t<1$, and one for
$t\geq1$. We insert a dead zone by setting
\begin{align*}
 q_a(t)&=q\left(\frac{t-a}{1-a}\right),\qquad a=\frac1{16},\\
 \psi(u)&=\frac1N\sum_{j=1}^N q_a\left(\frac N4u_j\right).
\end{align*}
With $\psi(z^0)=1$, define phase scores
\[
 \sigma_m(z)=\phi(z^m)+\psi(z^{m-1})-\frac54
\]
and the hard function
\begin{equation}
 \bar f(z)=\frac2N\left(\|[\sigma(z)]_+\|_2
 -\sum_{m=1}^M\psi(z^m)\right).
 \label{eq:base}
\end{equation}
This is the tilted-block construction of \citet{li2026deterministic} with
$q$ replaced by $q_a$.

The analytic structure is visible from a variational representation. Let
$W=\{w\in\R_+^M:\|w\|\leq1\}$ and set $w_{M+1}=0$. The support-function
identity for the positive orthant gives
\begin{align}
 \bar f(z)&=\max_{w\in W}F_w(z),\nonumber\\
 F_w(z)&=\frac2N\left(w_1-\frac54\sum_{m=1}^Mw_m\right)
 +\frac2N\sum_{m=1}^Mw_m\phi(z^m)\nonumber\\
 &\quad-\frac2N\sum_{m=1}^M(1-w_{m+1})\psi(z^m).
 \label{eq:variational}
\end{align}
Every $F_w$ has weak-convexity modulus at most $64/75$. Indeed, the tilted
terms are convex and the negative progress terms have this common curvature
bound. The pointwise maximum preserves it. The same formula gives the full
max-rule subdifferential used below, including all active tilted facets.

Let $L=\lfloor3N/16\rfloor$ and assign coordinate $(m,j)$ the level
\[
 \ell(m,j)=(m-1)L+j-1.
\]
Let $V_t$ span the coordinates of level at most $t$, with $V_{-1}=\{0\}$.
Every layer has at most six coordinates. Let $P_t$ denote orthogonal
projection onto $V_t$ and set $\tau=1/(64N^2)$.

\begin{proposition}[Robust chain certificates]
\label{prop:base}
The function $\bar f$ is bounded below, one-Lipschitz, and one-weakly convex.
It satisfies
$\bar f(0)-\inf\bar f\leq3M/N$. Moreover:
\begin{enumerate}
\item For any integer $t\geq-1$, if $|z_{m,j}|<\tau$ whenever
$\ell(m,j)>t$, then throughout a
neighborhood of $z$,
\[
 \bar f(y)=\bar f(P_{t+1}y).
\]
\item Set $b=q_a(3/4)=2783/3375$. If $\psi(z^M)<b$, then
$R(z):=\|[\sigma(z)]_+\|_2>0$. Every $g\in\partial\bar f(z)$ is
coordinatewise nonpositive and satisfies
\[
 \|g\|\geq\frac1{4\sqrt N}.
\]
If $g_{m,j}<0$ and $z^m_j>0$, then $z^m_j<13/(2N)$.
\end{enumerate}
\end{proposition}

The first statement is exact. For a partially revealed block, the first
hidden facet beats the next one by at least
\[
 \frac1{16N}-\frac{N\tau}{4}>0.
\]
For an unreached block, its phase score is at most
$-1/(16N)+N\tau/8<0$. All hidden progress coordinates lie strictly inside
the dead zone. To see the complete filtration argument, let
\[
 k_m=\min\{N,\max\{0,t-(m-1)L+1\}\}.
\]
If $1\leq k_m<N$, the progress term uses only coordinates through $k_m$, and
the strict facet comparison makes $k_m+1$ the only possibly relevant hidden
facet. If $k_m=0$ and $m\geq2$, then $k_{m-1}\leq L$, so
$\psi(z^{m-1})\leq3/16$ and the phase score is strictly negative. When
$m=1$ and $t=-1$, only its first facet can matter. Fully revealed blocks
already lie in $V_t$. The strict margins persist on a neighborhood, proving
local equality with the projection onto $V_{t+1}$.

For the second statement, $R(z)>0$ makes the active outer weight unique. The
tilted-facet and negative-progress terms in every subgradient are both
coordinatewise nonpositive, so they cannot cancel. Their block norms sum to
at least $1/(16N)$. Positivity of a supporting coordinate forces either an
active facet above $1/4$ or an unsaturated progress derivative. Both give the
$13/(2N)$ bound. Complete proofs appear in the supplement.

\section{A Lipschitz Wrapper for Unbounded Queries}
\label{sec:wrapper}

Let $n=MN$, let $U\in\R^{d\times n}$ satisfy $U^\top U=I_n$, and set
\begin{align*}
 B&=512\sqrt M,\qquad S=16B,\\
 s_S(x)&=\frac{x}{\sqrt{1+\|x\|^2/S^2}}.
\end{align*}
Define $H_U$ by \eqref{eq:wrapper}. Unlike the quadratic wrapper used for
smooth lower bounds, its radial term has slope at most three.

\begin{proposition}[Wrapper certificates]
\label{prop:wrapper}
For every $U$, the function $H_U$ is bounded below, one-Lipschitz, and
one-weakly convex. Its initial gap is at most $3M/(4N)$. If $x$ satisfies
\begin{equation}
 \left|\left\langle u_{m,j},s_S(x)\right\rangle\right|<\tau,
 \qquad j>\lfloor N/2\rfloor,
 \label{eq:terminal-hidden}
\end{equation}
then
\[
 \|\nabla(H_U)_{1/2}(x)\|>\frac1{64\sqrt N}.
\]
\end{proposition}

The class certificates follow from composition calculus. The map $s_S$ is
one-Lipschitz and its Jacobian is $6/S$-Lipschitz. Composing a one-Lipschitz,
one-weakly convex function with this map gives weak-convexity modulus at most
$1+6/S$. The factor $1/4$ restores modulus and Lipschitz constant below one.
The radial term is convex and nonnegative.

We outline the stationarity argument. Suppose the claimed bound fails, set
$z=\prox_{\frac12 H_U}(x):=\arg\min_y\{H_U(y)+\|y-x\|^2\}$, and write
$d=2(x-z)$. Then
$d\in\partial H_U(z)$. If $\|z\|\geq B$, the radial subgradient has norm at
least $3/\sqrt2$ before the factor $1/4$, while the hard term has norm at most
one. Hence $\|z\|<B$.

Let $\alpha=(1+\|z\|^2/S^2)^{-1/2}$ and $w=U^\top z$. For some
$g\in\partial\bar f(\alpha w)$, projection of $4d$ onto the columns of $U$
gives
\begin{equation}
 U^\top(4d)=\alpha g+c w,
 \qquad
 c=\frac3{\sqrt{B^2+\|z\|^2}}
 -\frac{\alpha^3\langle w,g\rangle}{S^2}.
 \label{eq:cancellation}
\end{equation}
Here $0<c<(3+1/256)/B$. Condition \eqref{eq:terminal-hidden}, the proximal
displacement, and the quadratic bound
$q_a(t)\leq(256/75)t_+^2$ imply
$\psi((\alpha w)^M)<b$. Proposition~\ref{prop:base} applies.

Only positive entries of $w$ can cancel the nonpositive vector $g$. On
$\supp(g)$, every such entry is at most $13/(2\alpha N)$. Therefore
\[
 \|\alpha g+cw\|
 \geq \left[
 \frac{\alpha}{4}
 -\frac{(3+1/256)13}{1024\alpha}
 \right]\frac1{\sqrt N}
 >\frac1{5\sqrt N}.
\]
After dividing by four, this contradicts
$\|d\|\leq1/(64\sqrt N)$.

\section{Random Rotation and Proof of the Main Theorem}
\label{sec:rotation}

Draw $U$ uniformly from the Stiefel manifold of $d\times n$ orthonormal
frames. Order its columns by nondecreasing filtration level, breaking ties in
a fixed way. We use a direct adapted conditional-Haar fact. If $v_k$ is
measurable with respect to an independent seed and columns
$u_1,\ldots,u_{k-1}$, with $\|v_k\|\leq S$, then conditional on this history
every remaining column has a uniform spherical marginal in a subspace of
dimension at least $d-n+1$. The standard spherical-cap bound used by
\citet[Appendix~B.3]{carmon2020stationary}, followed by a union bound, gives
$|\langle u_j,v_k\rangle|<\tau$ simultaneously for all $k\leq j$ when
\eqref{eq:dimension} holds.

\begin{lemma}[Robust random rotation]
\label{lem:rotation}
Fix a randomized algorithm and $\delta\in(0,1)$. If
\begin{equation}
 d\geq n+2S^2\tau^{-2}\log\frac{2n^2}{\delta},
 \label{eq:dimension}
\end{equation}
then, with probability at least $1-\delta$ over $U$ and the algorithm, every
query $x_t$ before filtration level $H$ satisfies
\[
 |\langle u_{m,j},s_S(x_t)\rangle|<\tau
 \quad\text{whenever }\ell(m,j)\geq t.
\]
\end{lemma}

Here is the reduction to that geometric fact. Let $\pi(1),\ldots,\pi(n)$ be
the column order and let $K_t=\{k:\ell(\pi(k))=t\}$. Define a surrogate
transcript on every history by replacing $\bar f(y)$ with $\bar f(P_ty)$ in
the wrapped oracle response at round $t$. Its effective query
$\widetilde y_t=s_S(\widetilde x_t)$ is measurable from the random seed and
columns in levels below $t$. For each $k\in K_t$, set
$v_k=\widetilde y_t$, and pad unused later entries with zero. Reordering a
Haar frame preserves its law. Moreover, $v_k$ does not use any column in its
own layer, so the sequence is adapted even when a query is repeated across a
width-six layer.

Apply the conditional-Haar fact to this virtual sequence. On its good event,
suppose the surrogate and true transcripts agree through query $t$. Every
coordinate of $U^\top s_S(x_t)$ at level at least $t$ then has magnitude below
$\tau$. Proposition~\ref{prop:base}, with index $t-1$, says that $\bar f$
agrees locally with $\bar f\circ P_t$. Thus their values and entire
subdifferentials agree. The exact chain rule through $U^\top s_S$ makes the
two wrapped oracle responses equal, and the next queries agree by induction.
This proves the lemma. The supplement provides the full adapted-Haar statement
and the measurability details.

We finish the construction. Let
\begin{align*}
 \kappa&=\frac1{64},\qquad
 r=\min\left\{\frac{\rho\Delta}{G^2},1\right\},\\
 N&=\left\lfloor\frac{\kappa^2G^2}{\ve^2}\right\rfloor,
 \qquad M=\left\lfloor\frac{Nr}{12}\right\rfloor.
\end{align*}
Under the accuracy assumption of Theorem~\ref{thm:main}, these integers are
admissible. Scale the wrapper as
\begin{equation}
 F_U(x)=\frac{G^2}{\rho}
 H_U\left(\frac{\rho}{G}x\right).
 \label{eq:scaling}
\end{equation}
Then $F_U\in\cF(G,\rho,\Delta)$ and
\[
 \|\nabla(F_U)_{1/(2\rho)}(x)\|
 =G\left\|\nabla(H_U)_{1/2}
 \left(\frac{\rho}{G}x\right)\right\|.
\]

Set $H=(M-1)L+\lfloor N/2\rfloor$. We have $H\geq MN/120$. Apply
Lemma~\ref{lem:rotation} to the induced algorithm whose wrapper query is
$\rho x/G$. Index its oracle calls as $x_0,\ldots,x_{T-1}$, and append its
transcript-measurable output as $x_T$. If $T<H$, the output round is still
before level $H$. Lemma~\ref{lem:rotation} and
Proposition~\ref{prop:wrapper} therefore show that the output is not
$\ve$-stationary with probability at least $1-\delta$. Set $\delta=1/4$.
Averaging over $U$ gives a fixed $U$ on which the algorithm succeeds with
probability at most $1/4$, which is strictly below the required $1/2$. Using
$N=\Omega(G^2/\ve^2)$ and $M=\Omega(Nr)$ proves
Theorem~\ref{thm:main}.

The dimension in \eqref{eq:dimension} is
$O(MN^4\log(MN))$, and hence polynomial in the inverse accuracy. We have not
optimized numerical constants.

\section{Discussion}
\label{sec:discussion}

Theorem~\ref{thm:main} removes randomization as a possible explanation for
the fourth-power complexity of weakly convex optimization. The obstruction
persists under a full-subdifferential oracle, so it does not arise from an
adversarial choice of one subgradient. In the small-gap regime, the existing
upper bound makes the nonsmooth penalty in \eqref{eq:minimax-rate} minimax
exact up to numerical constants.

Three qualifications are important. First, the construction uses a high but
polynomial dimension. Dimension-sensitive rates remain open. Second, the
oracle is local and first order. A proximal oracle can reveal global
information and supports faster rates for structured subclasses. Third, the
dead zone is essential for a full-set oracle. Random rotation alone gives only
small hidden projections, while any nonzero progress derivative would encode
an unrevealed column exactly.

The wrapper may be useful beyond this problem. Smooth lower bounds can use a
quadratic radial penalty because global function Lipschitzness is absent. The
linear-growth penalty in \eqref{eq:wrapper} retains the same information
shield while respecting a global Lipschitz budget. Its noncancellation proof
uses a sign-and-support certificate rather than smooth gradient geometry.

\bibliography{references}

\end{document}
